Hearing aids and hearables must enhance speech within about 5\,ms, so
they cannot use the future context that strong separators rely on. We ask
how much a strictly causal binaural enhancer loses as its algorithmic
latency drops from 32 to 1\,ms, and where that loss comes from. Real speech
and noise are spatialised through measured heads in simulated rooms. We
compare oracle masks, causal networks (an asymmetric-window GRU with deep
filtering over past frames, a symmetric-window twin, and a binaural
Conv-TasNet), their non-causal twins, and hop-matched lookahead. The window
costs little: an ideal mask gains \OracleEarAsymOnems\,dB SI-SDRi at 1\,ms
against \OracleEarAsymThreeTwoms\,dB at 32\,ms. The future costs almost
nothing: at 4\,ms, non-causal and causal twins differ by
\LongGapFutureFourms\,dB, and matched-hop lookahead changes quality by
under 0.05\,dB. Learning dominates: trained models remain
\LongGapLearningFourms\,dB below the oracle. Interaural level differences
stay within the 1\,dB JND, but no model, and not even the oracle, keeps the
ITD error under 20\,$\mu$s (best \stftasymFourmsearlongItd\,$\mu$s). A 4\,ms
causal model streams at real-time factor \RtfstftasymFourms\ on one CPU
thread. Short latencies are not what holds small causal enhancers back;
training and binaural timing fidelity are.
